% Generated from project outputs. Do not edit by hand.
\begin{table}[!htbp]
\centering
\begin{threeparttable}
\caption{Including in-kind by subsample: 2019 moments}
\label{tab:annual-2019-y2-moments}
\small
\setlength{\tabcolsep}{4pt}
\renewcommand{\arraystretch}{1.15}
\begin{tabularx}{\linewidth}{@{}Xrrrrr@{}}
\toprule
Sample & N & Mean & SD & Min & Max \\
\midrule
General & 94 & 2,781.2 & 1,577.9 & 376.0 & 6,242.2 \\
Men & 94 & 2,996.0 & 1,684.9 & 433.7 & 8,313.1 \\
Women & 94 & 2,423.3 & 1,683.5 & 254.0 & 7,672.2 \\
Urban & 94 & 2,978.0 & 1,612.7 & 317.3 & 6,389.7 \\
Rural & 94 & 2,370.0 & 1,512.8 & 333.6 & 9,510.9 \\
Indigenous & 94 & 2,170.7 & 1,424.6 & 237.4 & 7,537.6 \\
Non-Indigenous & 94 & 3,020.0 & 1,640.8 & 325.6 & 6,726.7 \\
Bottom 99 percent & 94 & 2,561.8 & 1,351.0 & 384.8 & 5,324.5 \\
Bottom 90 percent & 94 & 2,153.5 & 965.8 & 384.8 & 4,708.0 \\
Bottom 50 percent & 94 & 1,254.8 & 484.5 & 295.2 & 2,557.7 \\
\bottomrule
\end{tabularx}
\begin{tablenotes}[flushleft]
\footnotesize
\item Monthly quetzales. Unweighted distribution of survey-weighted cohort means before transition matching. N counts cohort cells, not individuals. SD is dispersion, not a standard error. Subsamples overlap. Cumulative income definitions and treatment of missing components follow the merging scripts.
\end{tablenotes}
\end{threeparttable}
\end{table}
